

\section{GPIOs} \label{GPIO}

%todo Warum hat es nicht funktioniert?
The pinout of the Jetson Nano and the instructions for installing the required libraries can be found in chapter \ref{Installation}.

\subsection{Control via the terminal}
The GPIOs can be controlled directly via the terminal. This is especially helpful for testing the functionality. For example, a circuit can be built with an LED that is controlled via board pin 12. A possible circuit is shown in the chapter \ref{LED}. 

\medskip

To get direct access to the GPIOs, you have to be a 'super user'.

\medskip
\SHELL{\$ sudo su}
\medskip

Now you can view the current configuration of the pins:

\medskip
\SHELL{\$ cat /sys/kernel/debug/gpio}
\medskip

To define pin 12 of the Jetson Nano as an output, tap

\medskip
\SHELL{\$ echo 79 > /sys/class/gpio/export} \newline
\SHELL{\$ echo out > /sys/class/gpio/gpio79/direction}
\medskip

Now the output of the pin can be set to high (3.3 V) or low (0 V):

\medskip
\SHELL{\$ echo 1 > /sys/class/gpio/gpio79/value} \newline
\SHELL{\$ echo 0 > /sys/class/gpio/gpio79/value}
\medskip

Finally, the pin should be released again:

\medskip
\SHELL{\$ echo 79 > /sys/class/gpio/unexport}
\medskip

\subsection{Verwendung der GPIO-Pins im Python-Programm}

If the GPIO pins are to be controlled from a Python programme, the following commands are elementary.

First, the library must be imported (before this, the steps described in the chapter \ref{Installation} for installing these libraries must be carried out). In order to be able to access the module under 'GPIO', the corresponding specification 'as GPIO' follows

\medskip
\SHELL{import Jetson.GPIO as GPIO}
\medskip

The Jetson GPIO library offers four options for numbering the I/O pins, so the numbering to be used must be specified.

\medskip
\SHELL{GPIO.setmode(GPIO.BOARD)} \newline
oder \newline
\SHELL{GPIO.setmode(GPIO.BCM)} \newline
oder \newline
\SHELL{GPIO.setmode(GPIO.CVM)} \newline
oder \newline
\SHELL{GPIO.setmode(GPIO.TEGRA\_SOC)}
\medskip

To set up a GPIO channel as an input or output, the commands

\medskip
\SHELL{GPIO.setup(channel, GPIO.IN)} 

or 

\SHELL{GPIO.setup(channel, GPIO.OUT)}
\medskip

where under \SHELL{channel} the pin is specified in the given numbering. The state of the pin prepared in this way can now be checked via

\medskip
\SHELL{GPIO.output(channel, GPIO.HIGH)}

and 

\SHELL{GPIO.output(channel, GPIO.LOW)}
\medskip

be changed.

\subsection{Visualisierung mit LED}\label{LED}

To visualise the circuitry of the GPIO pins, LEDs can be connected to the Jetson Nano. A possible circuit for this can be seen in figure \ref{Circuit}. It is explained in more detail under \cite{JetsonHacks:2020}. On the left side you can see the pins to which the elements are connected. Here, pin 2 has 5V, pin 6 is GND and pin 12 is switched to turn the LED on and off.

\begin{figure}[H]
	\begin{center}
%		\includegraphics[width=0.6\textwidth]{Jetson/Schaltung}
        
  \begin{tikzpicture}[scale = 0.75]
    \draw[line width=2pt] (0,0.5) -- (4,0.5) -- (4,1.5) -- (3.5,2) -- (4.0,2.5) -- (4,3.5) -- (4,5);
    \draw[line width=2pt] (4.0,5)  -- (4.5,5) -- (4.5,7)--(3.5,7)--(3.5,5)--(4,5);
    \draw[line width=2pt]  (3.5,3.5) -- (4.5,3.5);
    \draw[line width=2pt]  (4,3.5) -- (4.5,4.5) -- (3.5,4.5) -- cycle;
    \draw[-triangle 45]  (4.5, 4.2) -- (4.9,3.9);
    \draw[-triangle 45]  (4.45, 4.0) -- (4.85,3.7);
    \draw[line width=2pt] (4,7) -- (4,7.5) -- (0,7.5);
    \draw[line width=2pt] (0,2) -- (0.75,2) -- (0.75,2.5) -- (2.75,2.5) -- (2.75,1.5) -- (0.75,1.5) -- (0.75,2);
    \draw[line width=2pt] (2.75,2) -- (3.5,2);
    \draw[line width=2pt] ((3.5,1.5) -- (3.5,2.5);  
    \node[left] (Pin2) at (0,7.5)  {2};
    \node[left] (Pin12) at (0,2)  {12};
    \node[left] (Pin6) at (0,0.5)  {6};
    \node[above] (R10) at (1.75,2.5) {$10k\Omega$};
    \node[right] (BC) at (5,2) {BC337};
    \node[right] (LED) at (5,4) {LED};
    \node[right] (R330) at (5,6) {$330 \Omega$};
  \end{tikzpicture}        
		\caption{Circuit for testing the GPIO pins with LED}
		\label{Schaltung}
	\end{center}
\end{figure}

A photo of the circuit can be seen in figure \ref{CircuitPhoto}. The cable on the left in this photo is connected to pin 2, the one in the middle to pin 6 and the one on the right to pin 12.

\begin{figure}[H]
	\begin{center}
		\includegraphics[width=0.8\textwidth]{Jetson/SchaltungFoto}
		\caption{Circuit for testing the GPIO pins with LED}
		\label{SchaltungFoto}
	\end{center}
\end{figure}